\begin{statement}{B.14c}{Proved}{The displayed integral first cohomology groups vanish}
Put $T_{n,p}=\mathfrak s_{2,n}\otimes\mathbb Z_p$, with the
geometric integral $S_3$-action of B.10, and let
$\overline T_{n,p}=T_{n,p}/pT_{n,p}$ be the module called
$M_n$ in B.13. For $p=2,3$ and $1\leq n\leq9$,
\[
 H^1(S_3,T_{n,p})=0.
\]
Every positive-degree integral cohomology group is killed
by $p$. In the displayed range $H^2(S_3,T_{n,p})$ is an
$\mathbb F_p$-vector space of the following dimension:
\[
\begin{array}{c|rrrrrrrrr}
n&1&2&3&4&5&6&7&8&9\\\hline
p=2&0&0&1&0&1&0&2&0&3\\
p=3&0&0&0&0&0&0&0&1&0.
\end{array}
\]
Thus the nonzero modular first cohomology groups in B.13
do not give nonzero integral first cohomology obstructions
in these degrees.
\end{statement}
\begin{proof}
For $n\geq2$, B.10 gives the equivariant identification
\[
 T_{n,p}=\ker\bigl(V\otimes L_n(V)\longrightarrow L_{n+1}(V)\bigr),
 \qquad V=\mathbb Z_p^3/\mathbb Z_p(1,1,1).
\]
Here $x,y,z$ are the three images of the coordinate vectors,
$x+y+z=0$, and the indicated map is the bracket map.
The three peripheral conjugators are taken modulo their
common inner conjugator, giving precisely $V\otimes L_n$;
there is no further inner quotient for $n\geq2$.
The generator conventions are
$\rho(x)=y$, $\rho(y)=-x-y$, $\tau(x)=y$, $\tau(y)=x$.
These give exactly the tensor-action matrices used in B.13.
Changing the bracket map from $[v,a]$ to $[a,v]$ only
changes its global sign, not its kernel or its action.
In degree one the peripheral-centralizer quotient gives
$T_{1,p}=0$. Integral bracket surjectivity makes every
$T_{n,p}$ finite free and identifies its reduction with
the kernel calculated in B.13.

We record the coefficient argument explicitly. For a
finite group $G$ and an integral $G$-module $T$, the
homogeneous equivariant cochain complex admits the map
\[
 (hc)(g_0,\ldots,g_{q-1})
       =\sum_{s\in G}c(s,g_0,\ldots,g_{q-1}).
\]
It is equivariant by reindexing the sum, and direct
cancellation of the alternating-face terms gives
$dh+hd=|G|$ in positive degrees. Hence $|G|$ kills all
positive cohomology. For $G=S_3$ and $p=2,3$, its
prime-to-$p$ factor is a unit, so $p$ kills the integral
cohomology of $T_{n,p}$. These groups are finite-dimensional
$\mathbb F_p$-spaces: all cochain modules in each degree
are finite free over the Noetherian ring $\mathbb Z_p$.

Apply the coefficient sequence
$0\to T\xrightarrow{p}T\to T/pT\to0$ to the normalized
bar complex. It is a short exact sequence of cochain
complexes, since $S_3$ is finite and each cochain has
only finitely many values to lift. Its long exact sequence
contains
\[
 0\longrightarrow T^{S_3}/pT^{S_3}
 \longrightarrow(T/pT)^{S_3}
 \longrightarrow H^1(S_3,T)[p]\longrightarrow0
\]
and
\[
 0\longrightarrow H^1(S_3,T)/pH^1(S_3,T)
 \longrightarrow H^1(S_3,T/pT)
 \longrightarrow H^2(S_3,T)[p]\longrightarrow0.
\]

The rational invariant ranks from B.10, in degrees one
through nine, are
\[
 (0,0,1,0,1,0,2,1,4).
\]
The exact saved $d^0$ ranks in B.13 give this same list for
$\dim(\overline T_{n,p})^{S_3}$ at both primes: the latter
dimension is $\dim\overline T_{n,p}-\rank d^0$.
The script \texttt{code/g\_integral\_cohomology.py}
independently recomputes these modular invariant ranks
from the two saved generator matrices, verifies equality
with the saved $d^0$ ranks, and recomputes the rational
ranks by the character formula of B.10.
The first exact sequence therefore gives
$H^1(S_3,T_{n,p})[p]=0$. Since $p$ already kills the whole
group, its integral first cohomology vanishes.

The second exact sequence now identifies
\[
 H^1(S_3,\overline T_{n,p})
       \simeq H^2(S_3,T_{n,p})[p]
       =H^2(S_3,T_{n,p}).
\]
The $H^1$ column of B.13 gives the displayed dimensions.
All input and output ranks are recorded in
\texttt{g\_integral\_cohomology.json}.
This calculation by itself asserts no vanishing beyond
the displayed range and evaluates no specific higher
lifting obstruction.
\end{proof}
